% ========== TRUST & VERIFICATION ==========

\section{Trust \& Verification}
\label{sec:trust-verification}

\lettrine{T}{rust} and verification form the cornerstone of Symphony's extension security model, ensuring that extensions can be safely integrated into development workflows while maintaining system integrity and user confidence. The trust system implements multiple layers of verification, from cryptographic signatures to behavioral monitoring, creating a security framework.

\subsection{Code Signing Infrastructure}
\label{subsec:code-signing-infrastructure}

Symphony's code signing infrastructure provides cryptographic assurance of extension authenticity and integrity through a multi-tiered certificate authority system.

\subsubsection{Certificate Hierarchy}
\label{subsubsec:certificate-hierarchy}

The trust infrastructure employs a hierarchical certificate system that balances security with accessibility:

\begin{figure}[htbp]
\centering
\begin{tikzpicture}[
    cert/.style={rectangle, rounded corners, draw=brandPrimary, fill=brandPrimary!10, text width=3cm, text centered, minimum height=1cm},
    trust/.style={->, >=stealth, thick, brandSecondary}
]

% Certificate levels
\node[cert, fill=brandPrimary!30] (root) {Symphony Root CA};
\node[cert, below left=2cm and -1cm of root] (dev) {Developer CA};
\node[cert, below right=2cm and -1cm of root] (enterprise) {Enterprise CA};
\node[cert, below=2cm of dev] (community) {Community Certs};
\node[cert, below=2cm of enterprise] (internal) {Internal Certs};

% Trust relationships
\draw[trust] (root) -- (dev);
\draw[trust] (root) -- (enterprise);
\draw[trust] (dev) -- (community);
\draw[trust] (enterprise) -- (internal);

% Trust levels
\node[right=1cm of root, text width=4cm] {
    \textbf{Root Authority} \\
    \small Ultimate trust anchor \\
    \small Hardware security module
};

\node[right=1cm of dev, text width=4cm] {
    \textbf{Developer Certificates} \\
    \small Community developers \\
    \small Identity verification required
};

\node[right=1cm of enterprise, text width=4cm] {
    \textbf{Enterprise Certificates} \\
    \small Corporate extensions \\
    \small Enhanced validation
};

\end{tikzpicture}
\caption{Certificate Authority Hierarchy}
\label{fig:certificate-hierarchy}
\end{figure}

\subsubsection{Certificate Types and Trust Levels}
\label{subsubsec:certificate-types-trust-levels}

\begin{table}[h]
\centering
\caption{Certificate Types and Associated Trust Levels}
\label{tab:certificate-trust-levels}
\begin{tabular}{@{}llll@{}}
\toprule
\textbf{Certificate Type} & \textbf{Validation Level} & \textbf{Trust Level} & \textbf{Capabilities} \\
\midrule
Community & Identity verification & Standard & Public APIs only \\
Verified Developer & Enhanced validation & High & Extended APIs \\
Enterprise & Corporate validation & Very High & Internal APIs \\
Symphony Official & Internal validation & Maximum & System-level access \\
\bottomrule
\end{tabular}
\end{table}

\subsection{Signature Verification Process}
\label{subsec:signature-verification-process}

The signature verification process operates at multiple stages of the extension lifecycle, from initial installation through runtime execution.

\subsubsection{Multi-Stage Verification}
\label{subsubsec:multi-stage-verification}

\begin{figure}[htbp]
\centering
\begin{tikzpicture}[
    stage/.style={rectangle, rounded corners, draw=brandSecondary, fill=brandSecondary!10, text width=2.5cm, text centered, minimum height=1cm},
    check/.style={ellipse, draw=brandAccent, fill=brandAccent!10, text width=2cm, text centered, minimum height=0.8cm},
    flow/.style={->, >=stealth, thick, brandPrimary}
]

% Verification stages
\node[stage] (download) {Download};
\node[stage, right=2cm of download] (install) {Installation};
\node[stage, right=2cm of install] (load) {Loading};
\node[stage, right=2cm of load] (runtime) {Runtime};

% Verification checks
\node[check, below=1cm of download] {Integrity};
\node[check, below=1cm of install] {Signature};
\node[check, below=1cm of load] {Certificate};
\node[check, below=1cm of runtime] {Behavior};

% Flow arrows
\draw[flow] (download) -- (install);
\draw[flow] (install) -- (load);
\draw[flow] (load) -- (runtime);

% Verification arrows
\draw[flow, dashed] (download) -- (download |- 0,-1);
\draw[flow, dashed] (install) -- (install |- 0,-1);
\draw[flow, dashed] (load) -- (load |- 0,-1);
\draw[flow, dashed] (runtime) -- (runtime |- 0,-1);

\end{tikzpicture}
\caption{Multi-Stage Signature Verification Process}
\label{fig:signature-verification-stages}
\end{figure}

\subsubsection{Verification Algorithm}
\label{subsubsec:verification-algorithm}

The signature verification algorithm implements industry-standard cryptographic practices with Symphony-specific enhancements:

\begin{lstlisting}[language=rust,caption=Signature Verification Algorithm (Simplified)]
pub struct SignatureVerifier {
    root_ca: Certificate,
    trusted_cas: Vec<Certificate>,
    revocation_list: RevocationList,
}

impl SignatureVerifier {
    pub fn verify_extension(&self, extension: &Extension) -> VerificationResult {
        // 1. Extract signature and certificate chain
        let signature = extension.extract_signature()?;
        let cert_chain = extension.extract_certificate_chain()?;
        
        // 2. Verify certificate chain to root CA
        self.verify_certificate_chain(&cert_chain)?;
        
        // 3. Check certificate revocation status
        self.check_revocation_status(&cert_chain)?;
        
        // 4. Verify signature against extension content
        self.verify_signature(&signature, &extension.content, &cert_chain)?;
        
        // 5. Validate certificate permissions for extension type
        self.validate_permissions(&cert_chain, &extension.manifest)?;
        
        Ok(VerificationResult::Valid)
    }
}
\end{lstlisting}

\subsection{Trust Chains}
\label{subsec:trust-chains}

Trust chains establish transitive trust relationships that enable secure extension ecosystems while maintaining performance and usability.

\subsubsection{Chain of Trust Model}
\label{subsubsec:chain-of-trust-model}

Symphony implements a flexible chain of trust model that accommodates different organizational structures and security requirements:

\begin{infobox}[title=Trust Chain Principles]
\textbf{Transitivity}: Trust relationships can be inherited through verified chains \\
\textbf{Revocability}: Any link in the chain can be revoked, invalidating dependent trust \\
\textbf{Auditability}: All trust decisions are logged and auditable
\end{infobox}

\begin{table}[h]
\centering
\caption{Trust Chain Components}
\label{tab:trust-chain-components}
\begin{tabular}{@{}lll@{}}
\toprule
\textbf{Component} & \textbf{Role} & \textbf{Validation} \\
\midrule
Root Certificate & Ultimate trust anchor & Hardware-protected \\
Intermediate CA & Delegated authority & Regular audit required \\
Developer Certificate & Extension signer & Identity verification \\
Extension Signature & Content integrity & Cryptographic proof \\
Timestamp Authority & Signature timing & Independent validation \\
\bottomrule
\end{tabular}
\end{table}

\subsubsection{Trust Propagation}
\label{subsubsec:trust-propagation}

Trust propagates through the ecosystem based on established relationships and validation criteria:

\begin{compactlist}
\item \textbf{Direct Trust}: Extensions signed by directly trusted certificates
\item \textbf{Inherited Trust}: Extensions from developers trusted by verified organizations
\item \textbf{Conditional Trust}: Extensions with limited trust based on specific criteria
\item \textbf{Reputation-Based Trust}: Extensions with high community reputation scores
\end{compactlist}

\subsection{Revocation Mechanism}
\label{subsec:revocation-mechanism}

The revocation system provides rapid response capabilities for compromised certificates or malicious extensions while minimizing disruption to legitimate users.

\subsubsection{Revocation Infrastructure}
\label{subsubsec:revocation-infrastructure}

\begin{figure}[htbp]
\centering
\begin{tikzpicture}[
    component/.style={rectangle, rounded corners, draw=brandPrimary, fill=brandPrimary!10, text width=2.8cm, text centered, minimum height=1cm},
    data/.style={cylinder, draw=brandSecondary, fill=brandSecondary!10, text width=2cm, text centered, minimum height=1cm},
    flow/.style={->, >=stealth, thick, brandAccent}
]

% Components
\node[component] (crl) {Certificate Revocation List};
\node[component, right=3cm of crl] (ocsp) {OCSP Responder};
\node[data, below=2cm of crl] (database) {Revocation Database};
\node[component, below=2cm of ocsp] (monitor) {Behavior Monitor};

% Flows
\draw[flow] (database) -- (crl);
\draw[flow] (database) -- (ocsp);
\draw[flow] (monitor) -- (database);
\draw[flow, bend left=30] (crl) to node[above] {Updates} (ocsp);

% Labels
\node[above=0.5cm of crl] {\small Periodic Updates};
\node[above=0.5cm of ocsp] {\small Real-time Queries};
\node[left=0.5cm of database] {\small Central Store};
\node[right=0.5cm of monitor] {\small Detection};

\end{tikzpicture}
\caption{Certificate Revocation Infrastructure}
\label{fig:revocation-infrastructure}
\end{figure}

\subsubsection{Revocation Triggers}
\label{subsubsec:revocation-triggers}

Multiple events can trigger certificate or extension revocation:

\begin{expandedlist}
\item \textbf{Security Violations}: Malicious behavior, unauthorized access attempts
\item \textbf{Policy Violations}: Manifest misrepresentation, capability abuse
\item \textbf{Certificate Compromise}: Private key exposure, unauthorized usage
\item \textbf{Developer Misconduct}: Terms of service violations, malicious intent
\item \textbf{Technical Issues}: Severe bugs, system instability, data corruption
\end{expandedlist}

\subsubsection{Revocation Response}
\label{subsubsec:revocation-response}

The system implements graduated response mechanisms based on threat severity:

\begin{table}[h]
\centering
\caption{Revocation Response Levels}
\label{tab:revocation-response-levels}
\begin{tabular}{@{}llll@{}}
\toprule
\textbf{Severity} & \textbf{Response Time} & \textbf{Action} & \textbf{User Impact} \\
\midrule
Critical & Immediate & Block all instances & Immediate shutdown \\
High & 1 hour & Disable new installations & Existing continue with warnings \\
Medium & 24 hours & Mark as untrusted & User choice to continue \\
Low & 7 days & Add warnings & Minimal impact \\
\bottomrule
\end{tabular}
\end{table}

\subsection{User Consent Framework}
\label{subsec:user-consent-framework}

The user consent framework balances security with usability, providing clear information about extension capabilities while enabling informed decision-making.

\subsubsection{Consent Levels}
\label{subsubsec:consent-levels}

Symphony implements a tiered consent system that adapts to extension trust levels and capability requirements:

\begin{successbox}
\textbf{Smart Consent:} High-trust extensions with minimal capabilities require minimal user interaction, while untrusted extensions with extensive capabilities require explicit consent for each capability category.
\end{successbox}

\begin{table}[h]
\centering
\caption{User Consent Requirements}
\label{tab:user-consent-requirements}
\begin{tabular}{@{}llll@{}}
\toprule
\textbf{Trust Level} & \textbf{Capability Risk} & \textbf{Consent Required} & \textbf{User Experience} \\
\midrule
High & Low & Automatic & Silent installation \\
High & Medium & Notification & Install with notification \\
High & High & Confirmation & Single confirmation dialog \\
Medium & Any & Explicit & Detailed permission review \\
Low & Any & Granular & Per-capability consent \\
Untrusted & Any & Blocked & Installation prevented \\
\bottomrule
\end{tabular}
\end{table}

\subsubsection{Consent Interface Design}
\label{subsubsec:consent-interface-design}

The consent interface provides clear, actionable information about extension capabilities and risks:

\begin{compactlist}
\item \textbf{Risk Visualization}: Color-coded risk indicators and clear explanations
\item \textbf{Capability Grouping}: Related permissions grouped for easier understanding
\item \textbf{Impact Assessment}: Clear description of what each permission enables
\item \textbf{Alternative Options}: Suggestions for similar extensions with lower risk
\item \textbf{Granular Control}: Ability to grant subset of requested permissions
\end{compactlist}

\subsubsection{Ongoing Consent Management}
\label{subsubsec:ongoing-consent-management}

User consent is not a one-time decision but an ongoing relationship managed through the extension lifecycle:

\begin{alertbox}
\textbf{Dynamic Consent:} Extensions that exceed their declared capabilities or request additional permissions trigger new consent workflows, ensuring users maintain control over their security posture.
\end{alertbox}

The trust and verification system creates a robust security foundation that enables safe extension adoption while maintaining the flexibility and openness essential to Symphony's extension-first architecture.